\scanpage{92}
% loc 1
\[ \Delta(Q)=(bb')^2-4(b^2c'+b'^2c-4cc') \]
% loc 2
On trouve
% loc 3
\[ \Delta(Q)=(b^2-4c)(b'^2-4c')=\delta\delta' \tag*{(13)} \]
% loc 4–7
donc en fait, si \deleted{\unsure{(}} (comme on a supposé) $\delta\delta'\in k^{\unsure{\times|*}}$, alors $\unsure{Q(k)|Q(b)|\Delta(Q)}\in k^{\times}$ i.e ``racines distinctes''. Cela signifie aussi que \deleted{$t\ne bb'-t$ i.e $2t\ne bb'$}
% loc 8–10, cancelled lines
\[ \deleted{2\operatorname{Tr}(uv)=} \]
\[ \deleted{\delta\delta'\in k^{\unsure{1}*}\Longrightarrow 2t\ne bb'\ \text{partout}}\tag*{\deleted{(14)}} \]
\deleted{i.e $2t-bb'\in k^\times$}
% loc 11–14
\deleted{\unsure{Bref}} Le rev. des solutions de
\[ Q(t;b,c,b',c')\deleted{\unsure{t}}=0 \]
(en $t$, pour $b,c,b',c'$ fixés) est étale (quadratique) sur $S=\operatorname{Spec}k$.

% loc 15–19, left margin
% slanted
\margin{\deleted{\uline{A vérifier}. \unsure{Immédiat} sur corps, par ce qui précède}}
% loc 15–20, corollary
\uline{Cor.} (?) Soient $u,v$ dans alg. d'Azumaya \inserted{$A$}, tel que $\delta(u)\delta(v)\in k^*$, \deleted{\unsure{si|$k$}} et $u,v$ partout \inserted{lin. indép. mod 1.} \inserted{et que} $Q(u,v)=Q(t;b,c,b',c')=0$, \deleted{alors} il \deleted{$u,v$ appart. à une \unsure{même}} tels que $u,v$ appartiennent à une \deleted{\unsure{même}} sous-Alg. de Borel (qui n'est autre alors que $k\oplus ku\oplus kv$) $\Longleftrightarrow uv\in k.1+ku+kv$
% archive sheet number: 91
